\chapter{Mathematical Preliminaries}\label{ch02:chap}

This chapter assembles, at tutorial depth, every piece of mathematics the
kernel uses: finite sets and total functions; images, fibers, and the
classification of maps, with the counting facts that make bounded exhaustive
testing possible; the algebraic ladder from semigroups through monoids and
semigroupoids to categories; the structure theory of functional graphs of
endomaps, culminating in eventual periodicity and the factorial idempotent
power; the linear algebra of finite-dimensional Hilbert spaces, adjoints,
operator norms, and positive semidefinite matrices, including a complete
spectral computation for the all-ones matrix; and shortlex orders on words,
with the well-order property proved. All of it is standard. We cite it as
background --- to Hammack for sets and functions \cite{bookofproof}, to
Mac~Lane for category theory \cite{maclanecwm}, to Pin and to Howie for
finite semigroups and transformation monoids \cite{pin, howie}, and to Horn
and Johnson for matrix analysis \cite{horn} --- and we prove what we use, so
that the monograph is self-contained and so that later chapters can point at
exact statements rather than folklore. No novelty is claimed anywhere in
this chapter.

\section{Finite sets and total functions}\label{ch02:sec:finsets}

\subsection{Finite sets}

A set $X$ is \emph{finite} if it is empty or in bijection with
$\{0, 1, \dots, n-1\}$ for some positive integer $n$; in either case the
number $n \geq 0$ is unique and is written $|X|$, the \emph{cardinality} of
$X$. We take the elementary theory of finite sets as given
\cite{bookofproof} and record only the facts we lean on: a subset of a
finite set is finite with $|A| \leq |X|$, and $|A| = |X|$ for
$A \subseteq X$ forces $A = X$; a union of two finite sets is finite with
$|A \cup B| = |A| + |B| - |A \cap B|$; and the product of finite sets is
finite with $|X \times Y| = |X| \cdot |Y|$.

Throughout the monograph, finite sets come to us with more structure than
bare sethood: the kernel's objects carry an \emph{ordered} finite carrier, a
sequence $(x_0, x_1, \dots, x_{n-1})$ of pairwise distinct elements. The
order is data, not decoration --- Chapter~\ref{ch03:chap} is largely about
why --- but for the purposes of the present chapter, whenever we say
``finite set'' we permit ourselves to fix an enumeration and to identify the
set with its index set $\{0, \dots, n-1\}$ when convenient. Every
construction in this chapter is compatible with that identification, and the
kernel performs it explicitly via standard indices.

\subsection{Total functions}

\begin{definition}[Total function]\label{ch02:def:function}
A \emph{total function} $f : X \to Y$ between sets assigns to \emph{every}
element $x \in X$ exactly one element $f(x) \in Y$. The set $X$ is the
\emph{domain}, $Y$ the \emph{codomain}. Two total functions
$f, g : X \to Y$ with the same domain and codomain are \emph{equal} iff
$f(x) = g(x)$ for every $x \in X$; this is \emph{extensional} equality.
\end{definition}

Three points deserve emphasis, because each is load-bearing later.

First, \emph{totality is part of the definition}. A partial assignment ---
defined on some of $X$ --- is not a function $X \to Y$ in our usage. The
kernel enforces totality syntactically: an arrow's mapping tuple must supply
one target index for each of the $|X|$ source indices, and a tuple of the
wrong length is rejected at construction.

Second, \emph{the codomain is part of the data}. The same assignment rule
with a larger codomain is a different function; surjectivity, for instance,
is a property of the pair (rule, codomain), not of the rule alone. This
foreshadows the kernel's stricter stance, where even the codomain's
\emph{name} is part of an arrow's type.

Third, \emph{equality is extensional}. Two functions given by wildly
different descriptions --- one by a formula, one by a table, one as a
composite of seventeen other functions --- are equal the moment they agree
pointwise on a shared domain and codomain. The kernel mirrors this exactly:
arrow equality is equality of source, target, and mapping tuple, and the
presentation name attached to an arrow for diagnostics never participates
(Chapter~\ref{ch03:chap}).

For finite $X$ with a fixed enumeration $(x_0, \dots, x_{n-1})$ and finite
$Y$ with enumeration $(y_0, \dots, y_{m-1})$, a total function
$f : X \to Y$ is the same data as a tuple
$(t_0, \dots, t_{n-1}) \in \{0, \dots, m-1\}^n$ via
$f(x_i) = y_{t_i}$. This \emph{tuple representation} is the kernel's
concrete representation of arrows and will be used freely in proofs.

\subsection{Composition and identity}

\begin{definition}[Composition]\label{ch02:def:composition}
Given $f : X \to Y$ and $g : Y \to Z$, the \emph{composite}
$g \circ f : X \to Z$ is the total function $(g \circ f)(x) = g(f(x))$.
The \emph{identity} $\id_X : X \to X$ is the total function
$\id_X(x) = x$.
\end{definition}

\begin{proposition}[Associativity and neutrality]\label{ch02:prop:assoc}
For $f : X \to Y$, $g : Y \to Z$, $h : Z \to W$:
\[
h \circ (g \circ f) = (h \circ g) \circ f, \qquad
f \circ \id_X = f, \qquad \id_Y \circ f = f.
\]
\end{proposition}

\begin{proof}
All four sides of the first equation are total functions $X \to W$; for
each $x \in X$, both evaluate to $h(g(f(x)))$, so they are extensionally
equal. Similarly $(f \circ \id_X)(x) = f(\id_X(x)) = f(x)$ and
$(\id_Y \circ f)(x) = \id_Y(f(x)) = f(x)$ for all $x$.
\end{proof}

The proof is one line, and deliberately so: the entire content of the
category laws for total functions is pointwise evaluation plus extensional
equality. What the kernel adds in Chapter~4 is not mathematical depth but
\emph{typing}: the requirement that the middle objects of a composite agree
as declared objects, checked before the pointwise formula is ever consulted.

\section{Images, fibers, and the classification of maps}\label{ch02:sec:fibers}

\subsection{Image and fibers}

\begin{definition}[Image, fiber]\label{ch02:def:image-fiber}
For $f : X \to Y$: the \emph{image} is
$\im(f) = \{\, f(x) \mid x \in X \,\} \subseteq Y$; for $y \in Y$, the
\emph{fiber} over $y$ is $f^{-1}(y) = \{\, x \in X \mid f(x) = y \,\}
\subseteq X$.
\end{definition}

Fibers may be empty --- $f^{-1}(y) = \emptyset$ iff $y \notin \im(f)$ ---
and the notation $f^{-1}(y)$ does \emph{not} presuppose an inverse function.

\begin{lemma}[Fiber partition]\label{ch02:lem:fiber-partition}
The nonempty fibers of $f : X \to Y$ partition $X$: every $x \in X$ lies in
exactly one fiber, namely $f^{-1}(f(x))$. Consequently, for finite $X$,
\[
\sum_{y \in Y} \bigl| f^{-1}(y) \bigr| = |X|.
\]
\end{lemma}

\begin{proof}
By totality every $x$ has an image $f(x)$, so $x \in f^{-1}(f(x))$ and every
element lies in at least one fiber. If $x \in f^{-1}(y)$ and
$x \in f^{-1}(y')$ then $y = f(x) = y'$, so fibers over distinct points are
disjoint and every element lies in at most one. The counting identity is
then the statement that a partition's block sizes sum to the size of the
whole; empty fibers contribute zero and may be included in the sum harmlessly.
\end{proof}

This innocuous lemma does real work in Chapter~7: the trace identity
$\trace(A_f^{*} A_f) = |X|$ for the lift is precisely
Lemma~\ref{ch02:lem:fiber-partition} read off the diagonal of a Gram matrix,
and the norm theorem is a statement about the \emph{largest} block of the
fiber partition.

\subsection{Injections, surjections, bijections}

\begin{definition}\label{ch02:def:inj-surj}
A total function $f : X \to Y$ is \emph{injective} if $f(x) = f(x')$
implies $x = x'$; \emph{surjective} if $\im(f) = Y$; \emph{bijective} if it
is both.
\end{definition}

The fiber language rephrases all three: $f$ is injective iff every fiber
has at most one element; surjective iff every fiber is nonempty; bijective
iff every fiber has exactly one element. The kernel's isometry criterion
for lifts (Chapter~7) is the first of these in linear-algebraic clothing.

\begin{proposition}[Finite endomap trichotomy collapses]\label{ch02:prop:endo-trichotomy}
For a total function $f : X \to Y$ between finite sets with $|X| = |Y|$,
the following are equivalent: (i) $f$ is injective; (ii) $f$ is surjective;
(iii) $f$ is bijective.
\end{proposition}

\begin{proof}
(i) $\Rightarrow$ (ii): if $f$ is injective the nonempty fibers are
singletons, so by Lemma~\ref{ch02:lem:fiber-partition} the number of
nonempty fibers is $|X| = |Y|$; hence all $|Y|$ fibers are nonempty and $f$
is surjective. (ii) $\Rightarrow$ (i): if $f$ is surjective all $|Y|$
fibers are nonempty; since the fiber sizes are positive integers summing to
$|X| = |Y|$, each must equal $1$, so $f$ is injective. Either direction
gives (iii), and (iii) trivially gives (i) and (ii).
\end{proof}

\begin{remark}\label{ch02:rem:pigeonhole}
Proposition~\ref{ch02:prop:endo-trichotomy} is the pigeonhole principle in
functional dress, and it fails instantly for infinite sets (the successor
map on the natural numbers is injective, not surjective). Its relevance to
the kernel is indirect but real: on a finite carrier, a non-bijective
endomap must \emph{merge} somewhere and \emph{miss} somewhere, and repeated
iteration therefore shrinks the image until it stabilizes --- the
qualitative fact behind the eventual-periodicity theory of
Section~\ref{ch02:sec:funcgraphs}.
\end{remark}

\subsection{Composites and cancellation}

\begin{proposition}\label{ch02:prop:comp-inj-surj}
Let $f : X \to Y$ and $g : Y \to Z$. (i) If $f$ and $g$ are injective, so is
$g \circ f$; if both are surjective, so is $g \circ f$. (ii) If
$g \circ f$ is injective, $f$ is injective. (iii) If $g \circ f$ is
surjective, $g$ is surjective.
\end{proposition}

\begin{proof}
(i) If $g(f(x)) = g(f(x'))$, injectivity of $g$ gives $f(x) = f(x')$ and
injectivity of $f$ gives $x = x'$. For surjectivity, any $z \in Z$ is
$g(y)$ for some $y$, and $y = f(x)$ for some $x$, so $z = g(f(x))$.
(ii) If $f(x) = f(x')$ then $g(f(x)) = g(f(x'))$, so $x = x'$ by
injectivity of the composite. (iii) Any $z \in Z$ is $(g \circ f)(x)$ for
some $x$, hence $z = g(f(x)) \in \im(g)$.
\end{proof}

\subsection{Sections and retractions}

Injectivity and surjectivity admit a second characterization --- by
one-sided inverses --- that translates directly into the operator language
of Chapter~7, where it becomes the isometry and coisometry structure of
lifts.

\begin{proposition}[One-sided inverses]\label{ch02:prop:sections}
Let $f : X \to Y$ with $X$ nonempty and both sets finite.
\begin{enumerate}
\item $f$ is injective iff it has a \emph{retraction}: some
$r : Y \to X$ with $r \circ f = \id_X$.
\item $f$ is surjective iff it has a \emph{section}: some
$s : Y \to X$ with $f \circ s = \id_Y$.
\item $f$ is bijective iff it has a two-sided inverse, which is then
unique and equals every retraction and every section of $f$.
\end{enumerate}
\end{proposition}

\begin{proof}
(1) If $r \circ f = \id_X$ then $f$ is injective by
Proposition~\ref{ch02:prop:comp-inj-surj}(ii) applied to the injective
composite $\id_X$. Conversely, if $f$ is injective, fix any
$x_\ast \in X$ (here nonemptiness is used) and define $r(y) = x$ when
$y = f(x)$ for the (unique, by injectivity) such $x$, and
$r(y) = x_\ast$ when $y \notin \im(f)$. Then $r$ is total and
$r(f(x)) = x$ for every $x$.

(2) If $f \circ s = \id_Y$ then $f$ is surjective by
Proposition~\ref{ch02:prop:comp-inj-surj}(iii). Conversely, if $f$ is
surjective, every fiber $f^{-1}(y)$ is nonempty; choose for each $y$ an
element $s(y) \in f^{-1}(y)$ (a finite sequence of choices, no choice
axiom needed: pick the element of least standard index). Then
$f(s(y)) = y$ for every $y$.

(3) If $g$ satisfies $g \circ f = \id_X$ and $f \circ g = \id_Y$, then
$f$ is bijective by (1) and (2). Conversely a bijection's inverse
assignment $y \mapsto$ (the unique $x$ with $f(x) = y$) is a two-sided
inverse. For uniqueness and the coincidence claim: if
$r \circ f = \id_X$ and $f \circ s = \id_Y$, then
$r = r \circ \id_Y = r \circ (f \circ s) = (r \circ f) \circ s =
\id_X \circ s = s$; so any retraction equals any section, and in
particular two two-sided inverses are equal.
\end{proof}

\begin{remark}[Choice of section is data]\label{ch02:rem:section-choice}
The proof of (2) had to \emph{choose} an element of each fiber, and made
the choice canonical by taking least standard indices. This is a small
early instance of the chapter-spanning theme: where mathematics says
``choose arbitrarily,'' a deterministic kernel must say \emph{which}, and
the ordered carrier is what makes ``least index'' meaningful. An
implementation whose carriers were unordered sets could not even state
this canonical section.
\end{remark}

\section{Counting total maps}\label{ch02:sec:counting}

Bounded exhaustive testing --- the kernel's principal executable evidence
--- rests on knowing exactly how many objects of each shape exist, so that
``all of them'' is a checkable claim. This section proves the counts.

\begin{theorem}[Number of total maps]\label{ch02:thm:count-maps}
For finite sets $X$ and $Y$, the number of total functions $X \to Y$ is
$|Y|^{|X|}$, with the conventions $m^0 = 1$ (one empty function out of the
empty set, for any $Y$, including $Y = \emptyset$) and $0^n = 0$ for
$n \geq 1$ (no total function from a nonempty set into the empty set).
\end{theorem}

\begin{proof}
Write $n = |X|$, $m = |Y|$, and fix an enumeration
$x_0, \dots, x_{n-1}$ of $X$. We induct on $n$.

\emph{Base case $n = 0$.} A total function must assign an image to every
element of the domain; with no elements there is nothing to assign, and two
such functions are extensionally equal vacuously. So there is exactly one
function $\emptyset \to Y$ --- the empty function --- and $m^0 = 1$ counts
it correctly, even when $m = 0$.

\emph{Inductive step.} Suppose the count is $m^{n-1}$ for domains of size
$n - 1$, and let $|X| = n \geq 1$. If $m = 0$ there is no possible image
for $x_0$, hence no total function, and $0^n = 0$ is correct. If
$m \geq 1$, a total function $f : X \to Y$ is uniquely determined by the
pair (its restriction to $X \setminus \{x_0\}$, its value $f(x_0)$), and
every such pair arises from exactly one $f$: given a function $g$ on
$X \setminus \{x_0\}$ and a value $y \in Y$, define $f(x_0) = y$ and
$f(x_i) = g(x_i)$ for $i \geq 1$; this $f$ is total, restricts to $g$, and
takes the value $y$; and two distinct pairs yield extensionally distinct
functions. The correspondence is therefore a bijection between functions
$X \to Y$ and pairs, so the count is
$m^{n-1} \cdot m = m^n$ by the induction hypothesis and the product rule.
\end{proof}

\begin{corollary}[Counts used by the test families]\label{ch02:cor:test-counts}
(i) The number of endomaps of an $n$-element set is $n^n$; in particular
there are $3^3 = 27$ endomaps of a three-element carrier and hence
$27^3 = 19{,}683$ ordered triples of such endomaps --- the exhaustive
associativity family of Chapter~10. (ii) Summing over carrier sizes zero
through four, $\sum_{n=0}^{4} n^n = 1 + 1 + 4 + 27 + 256 = 289$ endomaps
--- the exhaustive endomap family (here $0^0 = 1$ counts the empty
function). (iii) The number of injective functions from an $n$-set to an
$m$-set is the falling factorial $m (m-1) \cdots (m-n+1)$ for
$n \leq m$ and $0$ otherwise; the number of bijections of an $n$-set with
itself is $n!$.
\end{corollary}

\begin{proof}
(i) and (ii) are immediate from Theorem~\ref{ch02:thm:count-maps}. For
(iii): choose images for $x_0, x_1, \dots$ in order; injectivity means each
choice must avoid the previously chosen images, leaving $m - i$ options at
step $i$; the product rule gives the falling factorial, which is $0$ as
soon as $n > m$ by pigeonhole
(Proposition~\ref{ch02:prop:endo-trichotomy}); bijections are the case
$n = m$.
\end{proof}

\begin{theorem}[Number of surjections]\label{ch02:thm:count-surj}
For finite sets with $|X| = n$ and $|Y| = m$, the number of surjective
total functions $X \to Y$ is
\[
\mathrm{Surj}(n, m) \;=\; \sum_{k=0}^{m} (-1)^{k} \binom{m}{k}
(m - k)^{n}.
\]
\end{theorem}

\begin{proof}
For each $y \in Y$ let $B_y$ denote the set of total functions
$X \to Y$ that \emph{miss} $y$, i.e.\ with $y \notin \im(f)$. The
surjections are the functions in none of the $B_y$, so by
inclusion-exclusion over the finite index set $Y$,
\[
\mathrm{Surj}(n, m)
= \sum_{S \subseteq Y} (-1)^{|S|} \Bigl| \bigcap_{y \in S} B_y \Bigr|.
\]
The intersection $\bigcap_{y \in S} B_y$ consists of the total functions
$X \to Y \setminus S$, of which there are $(m - |S|)^{n}$ by
Theorem~\ref{ch02:thm:count-maps}. Grouping the $2^m$ subsets by size
$k = |S|$, of which there are $\binom{m}{k}$ each, gives the stated sum.
\end{proof}

\begin{example}\label{ch02:ex:count-small}
On a three-element carrier: $27$ endomaps in all
(Theorem~\ref{ch02:thm:count-maps}); of these,
$\mathrm{Surj}(3,3) = 27 - 3 \cdot 8 + 3 \cdot 1 - 0 = 6$ are surjective
--- and indeed the surjective endomaps of a finite set are exactly its
bijections (Proposition~\ref{ch02:prop:endo-trichotomy}), of which there
are $3! = 6$; the two counts confirming each other is a useful sanity
identity, $\mathrm{Surj}(n,n) = n!$. The remaining $21$ endomaps of the
three-element carrier are the non-bijections, each of which merges at
least two points and misses at least one.
\end{example}

\begin{remark}[Exhaustibility, and where it stops]\label{ch02:rem:exhaustibility}
The counts explain the exact shape of the kernel's bounded exhaustive
test families (Chapter~10). Endomap families grow as $n^n$: sizes
$0$ through $4$ contribute $1 + 1 + 4 + 27 + 256 = 289$ cases ---
trivially exhaustible --- while size $5$ alone would add $3{,}125$ and
size $8$ would add $16{,}777{,}216$. Triples for associativity grow as
$(n^n)^3$: on the three-element carrier this is $27^3 = 19{,}683$, run
in milliseconds, while on a four-element carrier it would already be
$256^3 = 16{,}777{,}216$. The chosen bounds are the largest at which
\emph{complete} enumeration remains cheap enough to run on every test
invocation; beyond them the suite would have to sample, and sampled
evidence is a different, weaker kind of claim. Small-scope testing
doctrine \cite{jackson} holds that most implementation defects are
already expressible in small models; the kernel adopts that doctrine
explicitly and records its bounds rather than hiding them.
\end{remark}

\begin{remark}[Finiteness of the normal-form closure]\label{ch02:rem:closure-finite}
Theorem~\ref{ch02:thm:count-maps} is also the finiteness engine for
Chapter~8: the denotations reachable from a finite generator catalog over a
finite object set live inside the union of the sets of total maps between
catalog objects, of which there are at most $|Y|^{|X|}$ for each ordered
pair of objects, hence finitely many overall. The closure algorithm's
termination is exactly this observation.
\end{remark}

\section{The algebraic ladder: semigroups to categories}\label{ch02:sec:ladder}

The kernel's central type distinction --- endomaps have powers, paths do not
--- is a distinction between positions on a small ladder of algebraic
structures. We define each rung and prove the elementary facts; the reader
is referred to Howie \cite{howie} and Pin \cite{pin} for semigroup theory
and to Mac~Lane \cite{maclanecwm} for categories.

\begin{definition}[Semigroup, monoid]\label{ch02:def:semigroup}
A \emph{semigroup} is a set $S$ with an associative binary operation
$(s, t) \mapsto st$. A \emph{monoid} is a semigroup with a two-sided
identity element $1$: $1s = s1 = s$ for all $s$.
\end{definition}

\begin{proposition}[Uniqueness of identity; unambiguous powers]\label{ch02:prop:monoid-basics}
(i) A monoid's identity is unique. (ii) In a semigroup, the power $s^n$
($n \geq 1$), defined by $s^1 = s$ and $s^{n+1} = s^n s$, satisfies
$s^{m+n} = s^m s^n$ and $(s^m)^n = s^{mn}$; in a monoid the same laws hold
for all $m, n \geq 0$ with $s^0 = 1$.
\end{proposition}

\begin{proof}
(i) If $1$ and $1'$ are both identities, $1 = 1 \cdot 1' = 1'$. (ii)
$s^{m+n} = s^m s^n$ is induction on $n$: the case $n = 1$ is the
definition, and $s^{m + (n+1)} = s^{(m+n)+1} = s^{m+n} s =
(s^m s^n) s = s^m (s^n s) = s^m s^{n+1}$ using associativity. Then
$(s^m)^n = s^{mn}$ is induction on $n$ using the first law. The monoid
extension is immediate from $1 s = s$.
\end{proof}

The point of belaboring (ii): the expression $s^n$ is meaningful
\emph{because} all ways of parenthesizing $n$ copies of $s$ agree
(generalized associativity, itself a routine induction), and it presupposes
that the $n$ factors are \emph{the same element} of \emph{one} semigroup.
Both presuppositions fail for the generator sequences of
Section~\ref{ch01:sec:power-scenario}, and the kernel's refusal to give
paths a power operation is the refusal to notate a power whose defining
laws have no ambient semigroup in which to hold.

\begin{definition}[Category]\label{ch02:def:category}
A \emph{category} $\cC$ consists of a collection of \emph{objects}
$\Ob(\cC)$; for each ordered pair of objects $(X, Y)$ a set
$\Hom(X, Y)$ of \emph{arrows}; for each object an identity arrow
$\id_X \in \Hom(X, X)$; and a composition assigning to
$f \in \Hom(X, Y)$ and $g \in \Hom(Y, Z)$ an arrow
$g \circ f \in \Hom(X, Z)$, such that composition is associative wherever
defined and identities are neutral: $f \circ \id_X = f = \id_Y \circ f$
\cite{maclanecwm}.
\end{definition}

\begin{definition}[Semigroupoid]\label{ch02:def:semigroupoid}
A \emph{semigroupoid} is defined as a category is, but \emph{without} the
requirement of identity arrows: objects, hom-sets, and an associative
partial composition. Every category is a semigroupoid; a semigroupoid with
all identities is a category.
\end{definition}

\begin{proposition}[The ladder]\label{ch02:prop:ladder}
(i) A monoid is exactly a category with one object: the elements are the
arrows on that object, the operation is composition, the identity element
is the identity arrow. (ii) A semigroup is exactly a semigroupoid with one
object. (iii) In any category, each $\Hom(X, X)$ is a monoid under
composition, called the \emph{endomorphism monoid} $\mathrm{End}(X)$; in
any semigroupoid, each $\Hom(X, X)$ is a semigroup.
\end{proposition}

\begin{proof}
(i) Given a monoid $M$, define a category with a single object $\star$,
$\Hom(\star, \star) = M$, composition the monoid operation, and
$\id_\star = 1$; the category axioms are precisely associativity and the
identity laws. Conversely a one-object category's arrows form a monoid by
the same reading. (ii) is the same argument without identities. (iii)
Composition restricted to $\Hom(X, X)$ is total (source and target both
$X$), associative because the ambient composition is, and has identity
$\id_X$ (or, in a semigroupoid, is merely associative).
\end{proof}

\begin{remark}[Where powers live]\label{ch02:rem:where-powers-live}
Proposition~\ref{ch02:prop:ladder} locates the power operation precisely:
$f^n$ makes sense exactly when $f$ inhabits some $\Hom(X, X)$ --- a monoid
--- and makes no sense for $f \in \Hom(X, Y)$ with $X \neq Y$, since
already $f \circ f$ requires the target of the inner $f$ to be the source
of the outer. A composable sequence
$X_0 \to X_1 \to \cdots \to X_n$ of arrows in a category has a composite,
by iterated use of the composition operation and associativity, but the
sequence itself is an element of no single hom-set unless all the objects
coincide, and even then the \emph{sequence} is not an \emph{element}: it
composes to one, which is a different thing. The kernel's
\texttt{Endomap} / \texttt{SemigroupoidPath} type distinction (Chapters~5
and~6) is this remark made mechanical.
\end{remark}

\subsection{Generalized associativity}

Both the power notation $s^n$ and the composite of a semigroupoid path
rest on the same silent convention: an $n$-fold product may be written
without parentheses. That convention is a theorem, and since the kernel's
path composite (Chapter~6) is computed by one particular parenthesization
--- left-to-right folding --- while the mathematics quantifies over all
parenthesizations, we prove it.

\begin{theorem}[Generalized associativity]\label{ch02:thm:gen-assoc}
Let $a_1, a_2, \dots, a_n$ ($n \geq 1$) be a composable sequence in a
semigroupoid (in particular: any sequence in a semigroup, or a typed
path's arrow sequence). Then every way of fully parenthesizing the product
$a_n \cdots a_2 a_1$ yields the same element, which may therefore be
written without parentheses.
\end{theorem}

\begin{proof}
Define the \emph{left-folded} product $L_k$ by $L_1 = a_1$ and
$L_{k+1} = a_{k+1} L_k$. We show by strong induction on $n$ that every
fully parenthesized product $P$ of $a_n, \dots, a_1$ (in this order)
equals $L_n$. For $n = 1$, $P = a_1 = L_1$. For $n \geq 2$, the
outermost multiplication of $P$ splits the sequence at some position
$j$: $P = Q \cdot R$, where $R$ is a parenthesized product of
$a_j, \dots, a_1$ and $Q$ one of $a_n, \dots, a_{j+1}$, with
$1 \leq j \leq n - 1$. By induction $R = L_j$ and $Q = M$, where $M$ is
the left-folded product of $a_n, \dots, a_{j+1}$. If $j = n - 1$ then
$M = a_n$ and $P = a_n L_{n-1} = L_n$ directly. Otherwise write
$M = a_n M'$ with $M'$ the left-folded product of
$a_{n-1}, \dots, a_{j+1}$ (definition of left folding, since
$n - j \geq 2$). Using binary associativity,
\[
P = (a_n M') L_j = a_n (M' L_j),
\]
and $M' L_j$ is a parenthesized product of $a_{n-1}, \dots, a_1$, hence
equals $L_{n-1}$ by induction. So $P = a_n L_{n-1} = L_n$. (All
composites appearing are defined: in a semigroupoid, definedness of a
parenthesized product depends only on adjacent source-target matching of
the underlying sequence, which is given.)
\end{proof}

\begin{remark}\label{ch02:rem:fold-order}
Theorem~\ref{ch02:thm:gen-assoc} is what entitles the kernel to
\emph{implement} the path composite as a left fold while \emph{claiming}
the parenthesization-independent composite: the two coincide. It is also
the induction behind Proposition~\ref{ch02:prop:monoid-basics}(ii), made
explicit.
\end{remark}

\subsection{Structure-preserving maps: homomorphisms and functors}

One more definition from the standard toolkit is needed --- not for the
kernel's core, but for stating what Chapter~7 proves about the lift.

\begin{definition}[Monoid homomorphism; functor]\label{ch02:def:functor}
A \emph{monoid homomorphism} $\varphi : M \to N$ satisfies
$\varphi(1_M) = 1_N$ and $\varphi(st) = \varphi(s)\varphi(t)$. More
generally, a \emph{functor} $F : \cC \to \mathcal{D}$ between categories
assigns to each object $X$ of $\cC$ an object $F(X)$ of $\mathcal{D}$ and
to each arrow $f : X \to Y$ an arrow $F(f) : F(X) \to F(Y)$, such that
$F(\id_X) = \id_{F(X)}$ and $F(g \circ f) = F(g) \circ F(f)$ whenever
$g \circ f$ is defined \cite{maclanecwm}.
\end{definition}

By Proposition~\ref{ch02:prop:ladder}(i), a functor between one-object
categories is exactly a monoid homomorphism; the two clauses of
Definition~\ref{ch02:def:functor} are the two \emph{functor laws}.
Chapter~7's deterministic lift is an assignment
$X \mapsto H_X$, $f \mapsto A_f$ claimed to satisfy exactly these laws
($A_{\id_X} = I$ and $A_{g \circ f} = A_g A_f$); the claim's executable
verification --- exact matrix equality over all $1{,}678$ valid
composition pairs on carriers through size three --- and its partially
formalized status are detailed in Chapters~10 and~11. A functor
restricted to one endomorphism monoid is a monoid homomorphism
$\mathrm{End}(X) \to \mathrm{End}(H_X)$; in particular functors carry
powers to powers, $F(f^n) = F(f)^n$, by induction from the composition
law --- which is why the lift of an idempotent power is an idempotent
operator, a fact Chapter~7's counterexample then sharpens (idempotent,
yes; projector, no).

\subsection{The full transformation monoid}

\begin{definition}\label{ch02:def:full-transf}
For a finite set $X$, the \emph{full transformation monoid} $T_X$ is the
monoid of all total functions $X \to X$ under composition, with identity
$\id_X$ \cite{howie}.
\end{definition}

By Theorem~\ref{ch02:thm:count-maps}, $|T_X| = |X|^{|X|}$. The
invertible elements of $T_X$ are exactly the bijections
(Proposition~\ref{ch02:prop:endo-trichotomy} shows one-sided inverses
suffice in the finite case), forming the symmetric group of order $|X|!$.
Everything the kernel proves about endomap powers is a statement about
elements of $T_X$; the classical theory of finite transformation
semigroups \cite{howie, pin} contains far more (Green's relations,
ideal structure, generation), none of which the kernel needs or encodes.

\section{Functional graphs of endomaps}\label{ch02:sec:funcgraphs}

Fix a finite set $X$, $|X| = n \geq 1$, and an endomap $f : X \to X$. This
section develops the combinatorial structure of iteration --- the
\emph{functional graph} --- and proves the eventual-periodicity facts that
Chapter~5 packages as the safe finite power. The material is classical
\cite{pin, howie}.

\subsection{The rho shape}

\begin{definition}[Functional graph; orbit]\label{ch02:def:funcgraph}
The \emph{functional graph} of $f$ is the directed graph with vertex set
$X$ and one edge $x \to f(x)$ for each $x$; every vertex has out-degree
exactly one (totality and single-valuedness). The \emph{orbit} of $x$ is
the sequence $x, f(x), f^2(x), \dots$.
\end{definition}

\begin{definition}[Periodic and eventually periodic points]\label{ch02:def:periodic}
A point $x$ is \emph{periodic} if $f^k(x) = x$ for some $k \geq 1$; the
least such $k$ is its \emph{period}. A point is \emph{eventually periodic}
if $f^i(x)$ is periodic for some $i \geq 0$; the least such $i$ is the
\emph{tail length} (or index) of $x$.
\end{definition}

\begin{lemma}[Every orbit enters a cycle]\label{ch02:lem:rho}
Let $|X| = n$ and $x \in X$. Then:
\begin{enumerate}
\item The orbit of $x$ repeats within $n$ steps: there exist
$0 \leq i < j \leq n$ with $f^i(x) = f^j(x)$.
\item $x$ is eventually periodic, with tail length $t(x) \leq n - 1$ and
the period $\ell(x)$ of the periodic point it reaches satisfying
$1 \leq \ell(x) \leq n$, and moreover $t(x) + \ell(x) \leq n$.
\item For all $a \geq t(x)$ and all $b \geq 0$:
$f^{a + b}(x) = f^a(x)$ if and only if $\ell(x) \mid b$.
\end{enumerate}
\end{lemma}

\begin{proof}
(1) The $n + 1$ points $x, f(x), \dots, f^n(x)$ lie in the $n$-element set
$X$; by pigeonhole two coincide, say $f^i(x) = f^j(x)$ with
$0 \leq i < j \leq n$.

(2) Take $i < j$ as in (1) with $j$ minimal and, for that $j$, $i$
minimal. Set $t = i$ and $\ell = j - i$. The point $y = f^t(x)$
satisfies $f^\ell(y) = f^{j}(x) = f^{i}(x) = y$, so $y$ is periodic and
$x$ is eventually periodic with tail length at most $t$. Minimality of the
pair makes the points $x, f(x), \dots, f^{j-1}(x)$ pairwise distinct
(a repeat among them would contradict minimality of $j$, or of $i$ for the
given $j$); there are $t + \ell = j$ of these distinct points in an
$n$-set, so $t + \ell \leq n$, which gives both $t \leq n - 1$ (since
$\ell \geq 1$) and $\ell \leq n$ (since $t \geq 0$). Finally the period of
$y$ divides $\ell$ and its tail is $0$, so the stated bounds hold for the
true tail length and true period as well.

(3) Let $y = f^{t(x)}(x)$, a periodic point of period $\ell = \ell(x)$,
and let $a \geq t(x)$, say $a = t(x) + c$. Then $f^a(x) = f^c(y)$, and
$f^{a+b}(x) = f^{c + b}(y)$. Since $y$ has period $\ell$, the points
$y, f(y), \dots, f^{\ell - 1}(y)$ are pairwise distinct (a coincidence
$f^u(y) = f^v(y)$ with $0 \leq u < v < \ell$ would give
$f^{u + (\ell - v)}(y)= f^{\ell}(y) = y$ with
$1 \le u + \ell - v < \ell$, contradicting minimality of the period), and
$f^{m}(y) = f^{m'}(y)$ iff $m \equiv m' \pmod{\ell}$ for
$m, m' \geq 0$: the ``if'' direction is $f^{\ell}(y) = y$ iterated, and
the ``only if'' direction follows because both sides reduce to points of
the distinct list. Applying this with $m = c + b$, $m' = c$ gives
$f^{a+b}(x) = f^a(x)$ iff $\ell \mid b$.
\end{proof}

The picture painted by Lemma~\ref{ch02:lem:rho} is the classical
\emph{rho} ($\rho$) shape: each orbit runs along a tail of at most
$n - 1$ steps and then winds forever around a cycle of length at most $n$.
Globally, the functional graph decomposes into weakly connected components,
each consisting of exactly one directed cycle with rooted trees hanging off
the cycle vertices, edges oriented toward the cycle; the periodic points
are precisely the cycle vertices.

\subsection{Index and period of the endomap}

Iteration also acts on $f$ itself, inside the finite monoid $T_X$.

\begin{theorem}[Eventual periodicity of powers; index and period]\label{ch02:thm:index-period}
Let $f \in T_X$ with $|X| = n \geq 1$. There exist a least
$t \geq 0$ (the \emph{index}) and a least $p \geq 1$ (the \emph{period})
such that $f^{t + p} = f^{t}$; and then $f^{a} = f^{b}$ for
$a, b \geq t$ iff $a \equiv b \pmod p$. Moreover $t \leq n - 1$ and $p$
equals the least common multiple of the cycle lengths of the functional
graph, so $p \leq \mathrm{lcm}(1, 2, \dots, n)$ and, crudely,
$p \mid n!$.
\end{theorem}

\begin{proof}
The powers $f^0, f^1, f^2, \dots$ live in the finite monoid $T_X$ of size
$n^n$, so by pigeonhole some two coincide, $f^{i} = f^{j}$ with $i < j$;
then $f^{a} = f^{a + (j - i)}$ for all $a \geq i$ (compose both sides of
$f^i = f^j$ with $f^{a-i}$), so the sequence of powers is eventually
periodic as a sequence in $T_X$; define $t$ as the least index from which
some repetition recurs and $p$ as the least positive $b$ with
$f^{t+b} = f^t$. The biconditional for $a, b \geq t$ follows exactly as in
the proof of Lemma~\ref{ch02:lem:rho}(3), applied to the orbit of
$\id_X$ under composition with $f$ in $T_X$.

For the bound on $t$: by Lemma~\ref{ch02:lem:rho}(2) every point has tail
length at most $n - 1$, so for every $x$ the value $f^{n-1}(x)$ is already
periodic; hence for any $b$ that is a multiple of all point-periods,
$f^{(n-1) + b} = f^{n-1}$ pointwise, showing the index is at most
$n - 1$.

For the identification of $p$: let $L$ be the least common multiple of the
cycle lengths (equivalently, of the periods $\ell(x)$ over all $x$, since
each $\ell(x)$ is the length of the cycle its orbit enters). If
$a \geq t$ then in particular $a \geq t(x)$ for every $x$ (the pointwise
tails are dominated by the index --- indeed $f^a(x)$ must be periodic for
repetition of the map to occur at $a$), and
Lemma~\ref{ch02:lem:rho}(3) gives: $f^{a + b} = f^{a}$ as maps iff
$\ell(x) \mid b$ for every $x$, iff $L \mid b$. Hence $p = L$. Each cycle
length is at most $n$, so $L$ divides $\mathrm{lcm}(1, \dots, n)$, which in
turn divides $n!$ since every factor $1 \le k \le n$ divides $n!$.
\end{proof}

\subsection{The factorial idempotent power}

\begin{theorem}[Factorial idempotent power]\label{ch02:thm:factorial-idempotent}
For every finite set $X$ and every $f \in T_X$, the power $f^{|X|!}$ is
idempotent:
\[
f^{|X|!} \circ f^{|X|!} = f^{|X|!}.
\]
\end{theorem}

\begin{proof}
Write $N = |X|$. For $N = 0$ the only endomap is the empty function, which
is its own composite, and $f^{0!} = f^1 = f$; the claim holds. For
$N = 1$ the only endomap is $\id_X$, idempotent. Assume $N \geq 2$; then
$N! \geq N > N - 1$. Fix any $x \in X$. By Lemma~\ref{ch02:lem:rho}(2),
the tail length of $x$ is at most $N - 1 < N!$, so $f^{N!}(x)$ is a
periodic point; let $\ell \leq N$ be the length of its cycle. Since every
integer from $1$ to $N$ divides $N!$, we have $\ell \mid N!$, and
Lemma~\ref{ch02:lem:rho}(3) (with $a = N!$, $b = N!$) gives
$f^{2 N!}(x) = f^{N!}(x)$. As $x$ was arbitrary, $f^{2N!} = f^{N!}$
extensionally, which is the claim.
\end{proof}

\begin{remark}[Sharpness and cost]\label{ch02:rem:factorial-sharp}
The exponent $N!$ is chosen for the simplicity of its correctness proof,
not for minimality: Theorem~\ref{ch02:thm:index-period} shows the least
idempotent power of $f$ is $f^{a}$ for the least $a \geq t$ divisible by
$p$, which can be far smaller than $N!$. What $N!$ buys is
\emph{uniformity} --- one exponent that works for every endomap of every
$N$-element carrier, computable without inspecting the functional graph.
The kernel computes $f^{N!}$ by binary powering, so the cost is
$O(\log(N!)) = O(N \log N)$ compositions, each $O(N)$ array operations;
Chapter~5 details the implementation and the bounded presentation-name
scheme for such large exponents. The exhaustive verification of
Theorem~\ref{ch02:thm:factorial-idempotent} over all $289$ endomaps on
carriers of size at most four (Chapter~10) is possible precisely because
Corollary~\ref{ch02:cor:test-counts} counts that family exactly.
\end{remark}

\begin{remark}[What the theorem does not say]\label{ch02:rem:factorial-scope}
The theorem quantifies over single elements $f$ of the monoid $T_X$. It
says nothing about products of \emph{distinct} endomaps --- the alternating
scenario of Section~\ref{ch01:sec:power-scenario} satisfies the theorem for
each generator and still refutes the stabilization conclusion for the
sequence --- and nothing at all in the heterogeneous case, where no monoid
contains the sequence. This is the exact boundary the kernel's type system
draws.
\end{remark}

\subsection{Idempotents and the eventual image}

\begin{proposition}[Characterization of idempotents]\label{ch02:prop:idempotent-char}
For $e \in T_X$, the following are equivalent: (i) $e \circ e = e$;
(ii) $e$ restricts to the identity on its image, i.e.\ $e(y) = y$ for
every $y \in \im(e)$; (iii) $\im(e)$ is exactly the set of fixed points
of $e$.
\end{proposition}

\begin{proof}
(i) $\Rightarrow$ (ii): for $y \in \im(e)$ write $y = e(x)$; then
$e(y) = e(e(x)) = e(x) = y$. (ii) $\Rightarrow$ (iii): every image point
is fixed by (ii), and every fixed point $y = e(y)$ is an image point.
(iii) $\Rightarrow$ (i): for any $x$, $e(x) \in \im(e)$ is a fixed
point, so $e(e(x)) = e(x)$; extensionally $e \circ e = e$.
\end{proof}

\begin{proposition}[The eventual image]\label{ch02:prop:eventual-image}
Let $f \in T_X$, $|X| = n \geq 1$. The images of the powers form a
descending chain $\im(f) \supseteq \im(f^2) \supseteq \cdots$ that
stabilizes at some $k \leq n - 1$ to the \emph{eventual image}
$X_\infty = \im(f^{k}) = \im(f^{m})$ for all $m \geq k$. Moreover
$X_\infty$ is exactly the set of periodic points of $f$, $f$ restricts to
a bijection of $X_\infty$, and the idempotent $e = f^{n!}$ of
Theorem~\ref{ch02:thm:factorial-idempotent} has $\im(e) = X_\infty$.
\end{proposition}

\begin{proof}
$\im(f^{m+1}) = f(\im(f^m)) \subseteq \im(f^m)$ gives the chain. A
strictly decreasing chain of nonempty subsets of an $n$-set has length at
most $n$, and once two consecutive images coincide, all later ones do
(apply $f$ to the equality); since $|\im(f)| \le n$ and each strict step
drops the size by at least one, stabilization occurs by index $n - 1$ at
the latest (if no strict drop ever occurs, $f$ is bijective and the chain
is constant from the start).

Periodic points lie in every $\im(f^m)$ (a point of period $\ell$ is
$f^{m \ell}$ of itself), hence in $X_\infty$. Conversely, take
$m \geq n - 1$ and $x \in X$: the tail of $x$ has length at most
$n - 1 \leq m$ by Lemma~\ref{ch02:lem:rho}(2), so $f^{m}(x)$ lies on a
cycle and is periodic. Hence $\im(f^{m})$ consists of periodic points
for all such $m$, and $X_\infty$ is exactly the periodic set. On
$X_\infty$, $f$ is injective: a merge
$f(u) = f(v)$ with $u \neq v$ both periodic would give, for a common
multiple $P$ of their periods with $P \ge 1$,
$u = f^{P}(u)$ and $v = f^{P}(v)$, yet
$f^{P}(u) = f^{P-1}(f(u)) = f^{P-1}(f(v)) = f^{P}(v)$, forcing
$u = v$. An injective self-map of the finite set $X_\infty$
(note $f(X_\infty) = X_\infty$ by stabilization) is a bijection of it
(Proposition~\ref{ch02:prop:endo-trichotomy}). Finally
$\im(f^{n!}) = X_\infty$ since $n! \geq n - 1$ puts $n!$ beyond the
stabilization index.
\end{proof}

\begin{example}[A worked functional graph]\label{ch02:ex:funcgraph}
Let $X$ have six elements with standard indices $0, \dots, 5$ and let
$f$ have mapping tuple $(1, 2, 3, 1, 5, 5)$, i.e.
\[
0 \mapsto 1, \quad 1 \mapsto 2, \quad 2 \mapsto 3, \quad 3 \mapsto 1,
\quad 4 \mapsto 5, \quad 5 \mapsto 5.
\]
The functional graph has two components. Component one contains the
cycle $1 \to 2 \to 3 \to 1$ of length $3$, with the tail
$0 \to 1$ hanging on it: point $0$ has tail length $1$, points
$1, 2, 3$ are periodic of period $3$. Component two contains the fixed
point $5$ (a cycle of length $1$) with the tail $4 \to 5$: point $4$
has tail length $1$, point $5$ is periodic of period $1$.

Reading off the theory: tail lengths are $1$, comfortably within the
bound $n - 1 = 5$ of Lemma~\ref{ch02:lem:rho}; the cycle lengths are
$3$ and $1$, so the period of $f$ is $p = \mathrm{lcm}(3, 1) = 3$
(Theorem~\ref{ch02:thm:index-period}).

Compute the powers as mapping tuples, each a mechanical fold of
lookups:
\[
f = (1, 2, 3, 1, 5, 5), \qquad
f^{2} = (2, 3, 1, 2, 5, 5), \qquad
f^{3} = (3, 1, 2, 3, 5, 5), \qquad
f^{4} = (1, 2, 3, 1, 5, 5).
\]
Thus $f^{4} = f^{1}$ while $f^{1}, f^{2}, f^{3}$ are pairwise distinct
and none equals $f^{0} = \id_X$: the index of $f$ is $t = 1$ and the
period is confirmed as $p = 3$, with $f^{a} = f^{b}$ for
$a, b \geq 1$ iff $a \equiv b \pmod 3$.

The least idempotent power is the least $a \geq t$ divisible by $p$,
namely $a = 3$. Check directly against
Proposition~\ref{ch02:prop:idempotent-char}: $f^{3}$ sends
$0 \mapsto 3$, $1 \mapsto 1$, $2 \mapsto 2$, $3 \mapsto 3$,
$4 \mapsto 5$, $5 \mapsto 5$; its image is $\{1, 2, 3, 5\}$, which it
fixes pointwise --- an idempotent whose image is exactly the periodic
set, as Proposition~\ref{ch02:prop:eventual-image} requires. The
uniform exponent of Theorem~\ref{ch02:thm:factorial-idempotent} is
$6! = 720$; since $720 \geq 1$ and $720 \equiv 0 \pmod 3$, we get
$f^{720} = f^{3}$: the factorial power lands on the same idempotent, as
it must, at the cost of a much larger exponent --- the price of
uniformity noted in Remark~\ref{ch02:rem:factorial-sharp}.
\end{example}

The example also illustrates a methodological point the kernel takes
seriously: powers are computed and compared as mapping tuples, never
chased through orbits in prose. The tuple representation makes each
power an auditable mechanical object, and the exhaustive suite of
Chapter~10 compares exactly these tuples; informal orbit-chasing is
precisely where a wrong entry survives review.

\section{Finite-dimensional Hilbert spaces}\label{ch02:sec:hilbert}

The kernel's lift (Chapter~7) sends finite functions to linear maps between
finite-dimensional complex Hilbert spaces. We assemble the needed linear
algebra; the standard reference for this material is Horn and Johnson
\cite{horn}.

\subsection{Inner products and orthonormal bases}

\begin{definition}\label{ch02:def:hilbert}
A \emph{finite-dimensional (complex) Hilbert space} is a finite-dimensional
complex vector space $H$ with an inner product
$\langle \cdot, \cdot \rangle : H \times H \to \mathbb{C}$, linear in the
second argument, conjugate-linear in the first, conjugate-symmetric, and
positive definite ($\langle v, v \rangle > 0$ for $v \neq 0$). The norm is
$\lVert v \rVert = \sqrt{\langle v, v \rangle}$. A family
$(e_1, \dots, e_d)$ is an \emph{orthonormal basis} (ONB) if
$\langle e_i, e_j \rangle = \delta_{ij}$ and the family spans $H$.
\end{definition}

In finite dimension no completeness axiom is needed: every
finite-dimensional inner-product space is complete, so ``Hilbert space''
adds nothing beyond the inner product, and we use the term for continuity
with the lift literature. Every such space has an ONB (Gram--Schmidt), and
relative to an ONB every vector expands as
$v = \sum_i \langle e_i, v \rangle e_i$ with
$\lVert v \rVert^2 = \sum_i |\langle e_i, v \rangle|^2$ (Parseval; a
two-line computation from orthonormality).

\begin{proposition}[Cauchy--Schwarz]\label{ch02:prop:cauchy-schwarz}
For all $v, w \in H$: $|\langle v, w \rangle| \leq \lVert v \rVert
\lVert w \rVert$, with equality iff $v$ and $w$ are linearly dependent.
\end{proposition}

\begin{proof}
If $w = 0$ both sides vanish and $w = 0 \cdot v$ gives dependence.
Otherwise let $\lambda = \langle w, v \rangle / \lVert w \rVert^{2}$
and expand, using sesquilinearity and conjugate-symmetry:
\[
0 \;\leq\; \lVert v - \lambda w \rVert^{2}
\;=\; \lVert v \rVert^{2}
\;-\; \overline{\lambda}\,\langle w, v \rangle
\;-\; \lambda\,\overline{\langle w, v \rangle}
\;+\; |\lambda|^{2}\,\lVert w \rVert^{2}
\;=\; \lVert v \rVert^{2}
\;-\; \frac{|\langle w, v \rangle|^{2}}{\lVert w \rVert^{2}},
\]
because each of the three subtracted-or-added terms evaluates to
$\pm |\langle w, v \rangle|^{2} / \lVert w \rVert^{2}$ with signs
$-, -, +$. Rearranging gives
$|\langle w, v \rangle|^{2} \leq \lVert v \rVert^{2}
\lVert w \rVert^{2}$, and $|\langle v, w \rangle| =
|\langle w, v \rangle|$ by conjugate-symmetry. Equality holds iff
$\lVert v - \lambda w \rVert = 0$, i.e.\ $v = \lambda w$; conversely
dependence gives equality by direct substitution.
\end{proof}

Cauchy--Schwarz is used silently whenever the operator norm is
manipulated (for instance, it makes
$\lVert A \rVert = \max_{\lVert v \rVert = \lVert w \rVert = 1}
|\langle w, A v \rangle|$ an equivalent definition); we record it once
so that no later step is unproved.

For the kernel the relevant spaces are concrete: for a finite object $X$
with ordered carrier $(x_0, \dots, x_{n-1})$, the space $H_X$ is
$\mathbb{C}^n$ with the standard inner product and the standard ONB
$(e_{x_0}, \dots, e_{x_{n-1}})$, where the basis vector attached to a
carrier element is determined by that element's \emph{standard index}. The
order of the carrier is thereby welded to the geometry: permuting the
carrier permutes the basis, changing every matrix below. This is one of
the three places the ordered-input rule of Chapter~\ref{ch03:chap} is
load-bearing.

\subsection{Linear maps as matrices; adjoints}

Fix ONBs $(e_1, \dots, e_n)$ of $H_1$ and $(f_1, \dots, f_m)$ of $H_2$. A
linear map $A : H_1 \to H_2$ is determined by its values on the basis, and
its \emph{matrix} has entries $A_{ji} = \langle f_j, A e_i \rangle$, so
that $A e_i = \sum_j A_{ji} f_j$; composition of maps is multiplication of
matrices. The \emph{adjoint} $A^{*} : H_2 \to H_1$ is the unique linear
map with $\langle A^{*} w, v \rangle = \langle w, A v \rangle$ for all
$v \in H_1$, $w \in H_2$; existence and uniqueness in finite dimension
follow by defining $A^{*}$ on the basis via
$(A^{*})_{ij} = \overline{A_{ji}}$ --- the conjugate transpose --- and
checking the defining identity on basis pairs. The adjoint reverses
composition, $(AB)^{*} = B^{*} A^{*}$, and $(A^{*})^{*} = A$.

A map $U$ is an \emph{isometry} if $\lVert U v \rVert = \lVert v \rVert$
for all $v$; polarizing, this is equivalent to
$\langle Uv, Uw \rangle = \langle v, w \rangle$ for all $v, w$, and hence
to $U^{*} U = I$.

\subsection{Operator norm}

\begin{definition}\label{ch02:def:opnorm}
The \emph{operator norm} of a linear map $A : H_1 \to H_2$ between
finite-dimensional Hilbert spaces is
\[
\lVert A \rVert \;=\; \max_{\lVert v \rVert = 1} \lVert A v \rVert,
\]
the maximum being attained because the unit sphere of a finite-dimensional
space is compact and $v \mapsto \lVert A v \rVert$ is continuous.
\end{definition}

Immediately from the definition: $\lVert A v \rVert \leq \lVert A \rVert
\lVert v \rVert$ for all $v$; $\lVert A B \rVert \leq \lVert A \rVert
\lVert B \rVert$ (submultiplicativity); and $\lVert A \rVert = 0$ iff
$A = 0$.

\subsection{Positive semidefinite matrices and spectral facts}

\begin{definition}\label{ch02:def:psd}
A linear map $P : H \to H$ (equivalently its matrix in an ONB) is
\emph{self-adjoint} if $P^{*} = P$, and \emph{positive semidefinite} (PSD)
if it is self-adjoint and $\langle v, P v \rangle \geq 0$ for all $v$.
\end{definition}

For any $A : H_1 \to H_2$, the map $A^{*} A : H_1 \to H_1$ is PSD:
$(A^{*}A)^{*} = A^{*} A^{**} = A^{*} A$ and
$\langle v, A^{*} A v \rangle = \langle A v, A v \rangle =
\lVert A v \rVert^2 \geq 0$.

We use the spectral theorem for self-adjoint maps in the following form,
cited as standard \cite{horn}: \emph{every self-adjoint map $P$ on a
finite-dimensional Hilbert space has an ONB of eigenvectors with real
eigenvalues; $P$ is PSD iff all eigenvalues are nonnegative.} From it we
prove the two facts the lift theory consumes.

\begin{proposition}[Norm squared is the top eigenvalue of the Gram map]\label{ch02:prop:norm-eig}
For any linear $A : H_1 \to H_2$,
\[
\lVert A \rVert^{2} \;=\; \lambda_{\max}\bigl(A^{*} A\bigr),
\]
the largest eigenvalue of the PSD map $A^{*} A$.
\end{proposition}

\begin{proof}
Let $(u_1, \dots, u_n)$ be an ONB of $H_1$ consisting of eigenvectors of
$A^{*} A$, with eigenvalues $\lambda_1 \geq \cdots \geq \lambda_n \geq 0$.
For a unit vector $v = \sum_i c_i u_i$ (so $\sum_i |c_i|^2 = 1$),
\[
\lVert A v \rVert^{2}
= \langle A v, A v \rangle
= \langle v, A^{*} A v \rangle
= \Bigl\langle \sum_i c_i u_i, \sum_i c_i \lambda_i u_i \Bigr\rangle
= \sum_i \lambda_i |c_i|^{2},
\]
a convex combination of the eigenvalues, hence at most $\lambda_1$; and
the value $\lambda_1$ is attained at $v = u_1$. Taking the maximum over
unit $v$ and then square roots gives the claim.
\end{proof}

\begin{proposition}[Spectrum of the all-ones matrix]\label{ch02:prop:allones}
Let $J_m$ be the $m \times m$ matrix ($m \geq 1$) with every entry equal
to $1$. Then $J_m$ is PSD, $\rk(J_m) = 1$, $J_m^2 = m J_m$, and the
spectrum of $J_m$ is: eigenvalue $m$ with multiplicity $1$ (eigenvector
the all-ones vector $\mathbf{1} = (1, \dots, 1)^{\mathsf{T}}$), and
eigenvalue $0$ with multiplicity $m - 1$ (eigenspace the hyperplane
$\{ v : \sum_i v_i = 0 \}$). Consequently $\trace(J_m) = m$ and
$\lambda_{\max}(J_m) = m$.
\end{proposition}

\begin{proof}
$J_m = \mathbf{1} \mathbf{1}^{\mathsf{T}}$, an outer product, so every
column is a multiple of $\mathbf{1}$ and $\rk(J_m) = 1$; also
$J_m^{*} = J_m$ (real symmetric) and $\langle v, J_m v \rangle =
|\mathbf{1}^{\mathsf{T}} v|^{2} = \bigl| \sum_i v_i \bigr|^{2} \geq 0$, so
$J_m$ is PSD. Compute
$J_m \mathbf{1} = \mathbf{1} (\mathbf{1}^{\mathsf{T}} \mathbf{1}) =
m \mathbf{1}$: eigenvalue $m$ on the line spanned by $\mathbf{1}$. For
any $v$ with $\sum_i v_i = 0$,
$J_m v = \mathbf{1} (\mathbf{1}^{\mathsf{T}} v) = 0$: eigenvalue $0$ on a
hyperplane of dimension $m - 1$. The line and the hyperplane are
orthogonal complements, so together they exhaust $\mathbb{C}^{m}$, and
the listed spectrum is complete. Then
$J_m^{2} = \mathbf{1} (\mathbf{1}^{\mathsf{T}} \mathbf{1})
\mathbf{1}^{\mathsf{T}} = m J_m$, the trace is the eigenvalue sum
$m + 0 \cdot (m-1) = m$, and the largest eigenvalue is $m$.
\end{proof}

Two further pieces of standard machinery are needed to \emph{assemble}
spectra, because Chapter~7's Gram matrices arrive as block structures
only after a basis reordering.

\begin{lemma}[Unitary and permutation similarity preserve spectra]\label{ch02:lem:perm-similarity}
Let $U$ be unitary ($U^{*} U = U U^{*} = I$) and $P = U^{*} Q U$. Then
$P$ and $Q$ have the same eigenvalues with the same multiplicities, and
$P$ is PSD iff $Q$ is. In particular this applies when $U$ is a
\emph{permutation matrix} --- the matrix of a basis reordering, with
entries $U_{ij} = 1$ if the reordering sends position $j$ to position
$i$ and $0$ otherwise, which is unitary since its columns are distinct
standard basis vectors.
\end{lemma}

\begin{proof}
If $Q v = \mu v$ with $v \neq 0$, then
$P (U^{*} v) = U^{*} Q U U^{*} v = U^{*} Q v = \mu (U^{*} v)$ and
$U^{*} v \neq 0$ ($U^{*}$ is invertible), so $\mu$ is an eigenvalue of
$P$; the map $v \mapsto U^{*} v$ is a linear bijection between the
eigenspaces, preserving multiplicities, and the argument reverses via
$Q = U P U^{*}$. For positivity:
$\langle v, P v \rangle = \langle U v, Q\, U v \rangle$, and
$v \mapsto U v$ ranges over all vectors, so the two quadratic forms
are nonnegative together; self-adjointness transfers by
$P^{*} = U^{*} Q^{*} U$.
\end{proof}

\begin{lemma}[Spectrum of a direct sum]\label{ch02:lem:direct-sum}
Let $P = \bigoplus_{r=1}^{s} P_r$ be block-diagonal with square blocks
$P_1, \dots, P_s$. Then the spectrum of $P$ is the multiset union of
the spectra of the blocks, and $\lambda_{\max}(P) = \max_r
\lambda_{\max}(P_r)$ when all blocks are self-adjoint.
\end{lemma}

\begin{proof}
The ambient space decomposes as an orthogonal direct sum
$H = \bigoplus_r H_r$ with $P$ mapping each $H_r$ to itself by $P_r$.
A vector $v = (v_1, \dots, v_s)$ satisfies $P v = \mu v$ iff
$P_r v_r = \mu v_r$ for every $r$; for $v \ne 0$ some $v_r \neq 0$, so
every eigenvalue of $P$ is an eigenvalue of some block, and conversely a
block eigenvector extends by zeros to an eigenvector of $P$. For
multiplicities: the $\mu$-eigenspace of $P$ is the direct sum of the
$\mu$-eigenspaces of the blocks, so dimensions add. The
$\lambda_{\max}$ statement follows since the multiset union's maximum is
the maximum of the blockwise maxima.
\end{proof}

\subsection{Projectors and partial isometries}

Chapter~7's permanent counterexample --- an idempotent lift that is
\emph{not} a projector --- needs the following standard vocabulary to be
stated precisely.

\begin{definition}\label{ch02:def:projector}
A linear map $P : H \to H$ is an \emph{orthogonal projector} if
$P^{2} = P$ and $P^{*} = P$. A linear map $A : H_1 \to H_2$ is a
\emph{partial isometry} if $A^{*} A$ is an orthogonal projector.
\end{definition}

\begin{proposition}[Elementary projector facts]\label{ch02:prop:projector-facts}
(i) An orthogonal projector satisfies $\lVert P \rVert \leq 1$, with
$\lVert P \rVert = 1$ unless $P = 0$. (ii) An idempotent that is not
self-adjoint can have norm strictly greater than $1$; idempotence alone
bounds nothing. (iii) An isometry is a partial isometry
($A^{*} A = I$ is a projector).
\end{proposition}

\begin{proof}
(i) $P$ is self-adjoint, so by the spectral theorem it has an ONB of
eigenvectors; $P^2 = P$ forces each eigenvalue $\mu$ to satisfy
$\mu^{2} = \mu$, so $\mu \in \{0, 1\}$. By
Proposition~\ref{ch02:prop:norm-eig} applied to
$P^{*} P = P^{2} = P$, $\lVert P \rVert^{2} = \lambda_{\max}(P)
\leq 1$, with equality iff the eigenvalue $1$ occurs, i.e.\ iff
$P \neq 0$. (ii) is witnessed by any concrete example; the canonical
one, the $2 \times 2$ matrix with rows $(1, 1)$ and $(0, 0)$, is
exactly the lifted constant map of Chapter~7, where its norm is computed
to be $\sqrt{2}$. (iii) is immediate from the definitions since $I$ is
a projector.
\end{proof}

\begin{remark}[Why these two facts]\label{ch02:rem:why-these}
Chapter~7 shows that for the lift $A_f$ of a finite function $f$, the Gram
map $A_f^{*} A_f$ is, after grouping the input basis by fibers, a direct
sum of blocks $J_{m_y}$ with $m_y = |f^{-1}(y)|$ the fiber sizes. The
grouping is a basis reordering --- a permutation similarity, harmless to
spectra by Lemma~\ref{ch02:lem:perm-similarity} --- and the direct sum's
spectrum is the multiset union of the block spectra by
Lemma~\ref{ch02:lem:direct-sum}. Feeding
Proposition~\ref{ch02:prop:allones} through
Proposition~\ref{ch02:prop:norm-eig} then yields the exact norm formula
$\lVert A_f \rVert^{2} = \max_y |f^{-1}(y)|$ with no analysis beyond what
this section contains. The trace identity
$\trace(A_f^{*} A_f) = \sum_y m_y = |X|$ is
Lemma~\ref{ch02:lem:fiber-partition}. The primer is arranged so that
Chapter~7's proofs are assembly, not invention.
\end{remark}

\section{Words and the shortlex order}\label{ch02:sec:shortlex}

The kernel's normal forms (Chapter~8) select, for each reachable
denotation, the least word denoting it under the \emph{shortlex} order.
For that selection to be well defined and stable, shortlex must be a total
order with no infinite descent --- a well-order --- on the relevant word
sets. This section defines the order and proves those properties.

\subsection{Words over a ranked alphabet}

\begin{definition}[Words]\label{ch02:def:words}
Let $G$ be a finite set of \emph{letters} equipped with a fixed total
order, presented as a \emph{ranking}: a bijection
$\rho : G \to \{0, 1, \dots, |G| - 1\}$. A \emph{word} over $G$ is a
finite (possibly empty) sequence $w = g_{1} g_{2} \cdots g_{k}$ of
letters; $|w| = k$ is its \emph{length}, and $\varepsilon$ denotes the
empty word. $G^{*}$ is the set of all words.
\end{definition}

In the kernel the letters are generator labels from a declared catalog,
and the ranking is the catalog's declared order --- explicitly supplied,
never inferred from container iteration (Chapter~\ref{ch03:chap}). Words
are additionally \emph{typed} (their letters' arrows must compose in
sequence from a declared start object), which restricts $G^{*}$ to a
subset; every property proved below is inherited by subsets, so we work
with $G^{*}$ and note the inheritance at the end.

\begin{definition}[Shortlex]\label{ch02:def:shortlex}
For words $u = u_1 \cdots u_k$ and $v = v_1 \cdots v_l$ over $(G, \rho)$,
say $u <_{\mathrm{slex}} v$ iff either
\begin{enumerate}
\item $k < l$ (shorter words come first); or
\item $k = l$ and there is a position $i \leq k$ with
$\rho(u_j) = \rho(v_j)$ for all $j < i$ and $\rho(u_i) < \rho(v_i)$
(equal lengths compare lexicographically by rank).
\end{enumerate}
\end{definition}

\begin{proposition}[Shortlex is a strict total order]\label{ch02:prop:slex-total}
On $G^{*}$, the relation $<_{\mathrm{slex}}$ is irreflexive, transitive,
and total: for distinct words $u \neq v$, exactly one of
$u <_{\mathrm{slex}} v$, $v <_{\mathrm{slex}} u$ holds.
\end{proposition}

\begin{proof}
\emph{Irreflexivity}: $u <_{\mathrm{slex}} u$ would require $|u| < |u|$ or
a position where $\rho(u_i) < \rho(u_i)$; both are impossible.

\emph{Totality}: let $u \neq v$. If $|u| \neq |v|$ the length clause
orders them one way and (by irreflexivity of $<$ on integers) not the
other. If $|u| = |v| = k$, then since $u \neq v$ there is a least
position $i$ with $u_i \neq v_i$; all earlier positions agree, and since
$\rho$ is injective, exactly one of $\rho(u_i) < \rho(v_i)$,
$\rho(v_i) < \rho(u_i)$ holds, ordering the words one way and not the
other.

\emph{Transitivity}: suppose $u <_{\mathrm{slex}} v <_{\mathrm{slex}} w$.
If $|u| < |v|$ or $|v| < |w|$ then $|u| < |w|$ (lengths are ordered
integers and are nondecreasing along $<_{\mathrm{slex}}$), giving
$u <_{\mathrm{slex}} w$ unless $|u| = |w|$, which can happen only when
$|u| = |v| = |w|$. So assume all lengths equal $k$ and let $i$ (resp.\
$j$) be the deciding position for $u$ against $v$ (resp.\ $v$ against
$w$). Put $h = \min(i, j)$. For all positions $< h$, all three words
agree. At position $h$: if $h = i < j$ then $u_h$ has smaller rank than
$v_h = w_h$; if $h = j < i$ then $u_h = v_h$ has smaller rank than $w_h$;
if $h = i = j$ then $\rho(u_h) < \rho(v_h) < \rho(w_h)$. In every case
$h$ witnesses $u <_{\mathrm{slex}} w$.
\end{proof}

\subsection{The well-order property}

\begin{theorem}[Shortlex is a well-order]\label{ch02:thm:slex-wellorder}
Every nonempty subset $S \subseteq G^{*}$ has a
$<_{\mathrm{slex}}$-least element. Equivalently, there is no infinite
strictly $<_{\mathrm{slex}}$-descending sequence of words.
\end{theorem}

\begin{proof}
The key counting observation: for each fixed length $k$ there are exactly
$|G|^{k}$ words of length $k$ (Theorem~\ref{ch02:thm:count-maps}, reading
a word as a function from positions to letters), a finite number; and
every word of length $k$ precedes every word of length $k + 1$ in
shortlex. Hence for any word $w$, the set
$\{\, u \in G^{*} \mid u <_{\mathrm{slex}} w \,\}$ is contained in the
words of length at most $|w|$, of which there are
$\sum_{k=0}^{|w|} |G|^{k}$ --- finitely many. (For $|G| = 0$ the only word
is $\varepsilon$ and everything below is trivial; assume $|G| \ge 1$.)

Now let $S$ be nonempty and pick any $w_0 \in S$. The set
$S' = \{\, u \in S \mid u \leq_{\mathrm{slex}} w_0 \,\}$ is nonempty
(contains $w_0$) and finite (contained in the finite set of words at most
$w_0$, plus $w_0$ itself). A nonempty finite subset of a strict total
order has a least element --- induct on the size: remove any element; the
remainder's least element (or, at size one, the sole element) compared
with the removed one yields the least, using totality
(Proposition~\ref{ch02:prop:slex-total}). Let $m$ be the least element of
$S'$. For any $s \in S$: if $s \leq_{\mathrm{slex}} w_0$ then
$s \in S'$ and $m \leq_{\mathrm{slex}} s$; otherwise
$m \leq_{\mathrm{slex}} w_0 <_{\mathrm{slex}} s$ by transitivity. So $m$
is least in $S$.

For the equivalence with no infinite descent: an infinite strictly
descending sequence would be a nonempty set with no least element (its
every member is exceeded from below by a later member); conversely if some
nonempty $S$ had no least element one could choose a strictly descending
sequence by repeatedly picking a smaller member. We only use the first
direction, which is constructive given the finiteness above.
\end{proof}

\begin{corollary}[Minimal representatives exist and are unique]\label{ch02:cor:min-rep}
Let $W \subseteq G^{*}$ be any set of words and let $\sim$ be any
equivalence on $W$ (in Chapter~8: ``denotes the same finite function'').
Every nonempty equivalence class $C$ has a unique
$<_{\mathrm{slex}}$-least element, which may serve as the canonical
representative $\NF(C)$.
\end{corollary}

\begin{proof}
Existence is Theorem~\ref{ch02:thm:slex-wellorder} applied to
$C \subseteq G^{*}$. Uniqueness: two least elements are mutually
$\leq_{\mathrm{slex}}$, and antisymmetry --- immediate from totality and
irreflexivity of the strict order --- forces them equal.
\end{proof}

\subsection{Enumeration and the shortlex rank}

Shortlex is not only a well-order; it is a computable enumeration, and
the position of any word in that enumeration has a closed form. This
matters operationally: Chapter~8's closure explores words in exactly this
enumeration order, and the formula below is what makes ``explored in
shortlex order'' an auditable claim rather than an implementation hope.

\begin{proposition}[Shortlex rank formula]\label{ch02:prop:slex-rank}
Let $|G| = g \geq 1$ and let $w = w_1 \cdots w_k$ be a word of length
$k$. The number of words strictly $<_{\mathrm{slex}}$-below $w$ is
\[
\rk_{\mathrm{slex}}(w) \;=\; \sum_{j=0}^{k-1} g^{\,j}
\;+\; \sum_{i=1}^{k} \rho(w_i)\, g^{\,k-i},
\]
the first sum counting all strictly shorter words and the second the
length-$k$ words that are lexicographically smaller. The map
$w \mapsto \rk_{\mathrm{slex}}(w)$ is the unique order-isomorphism from
$(G^{*}, <_{\mathrm{slex}})$ onto $(\mathbb{N}, <)$.
\end{proposition}

\begin{proof}
Words shorter than $w$: there are $g^{j}$ words of each length
$j < k$ (Theorem~\ref{ch02:thm:count-maps}), all below $w$ by the
length clause, giving the first sum. Words of length exactly $k$ below
$w$: such a word $u$ agrees with $w$ on some (possibly empty) prefix of
length $i - 1$ and has $\rho(u_i) < \rho(w_i)$ at the first
disagreement $i$; for each choice of $i \in \{1, \dots, k\}$ there are
$\rho(w_i)$ choices of the smaller letter and $g^{\,k-i}$ free choices
for the remaining positions $i+1, \dots, k$, and these cases are
disjoint (distinct first-disagreement positions) and exhaustive. Summing
gives the second sum.

For the order-isomorphism claim: $\rk_{\mathrm{slex}}$ is strictly
monotone (if $u <_{\mathrm{slex}} w$ then every word below $u$ is below
$w$ and $u$ itself is below $w$, so
$\rk_{\mathrm{slex}}(u) < \rk_{\mathrm{slex}}(w)$), hence injective;
and it is surjective onto $\mathbb{N}$ because the count of words below
the $(n+1)$-st word of the enumeration is exactly $n$ (each natural
number is achieved as one steps through the enumeration, which exists by
Theorem~\ref{ch02:thm:slex-wellorder}: the least word has rank $0$, and
the least word above any $w$ has rank $\rk_{\mathrm{slex}}(w) + 1$
since finitely many words lie below any word). Uniqueness of an
order-isomorphism between two linear orders with no greatest element in
which every element has finitely many predecessors is standard: any two
such isomorphisms agree by induction along the enumeration.
\end{proof}

\begin{example}[Shortlex over a two-letter catalog]\label{ch02:ex:slex-two}
Let $G = \{ a, b \}$ with declared ranking $\rho(a) = 0$,
$\rho(b) = 1$. The shortlex enumeration begins
\[
\varepsilon,\; a,\; b,\; aa,\; ab,\; ba,\; bb,\; aaa,\; aab,\; aba,\;
abb,\; baa,\; \dots
\]
and, for instance, $\rk_{\mathrm{slex}}(ba) = (1 + 2) +
(1 \cdot 2 + 0 \cdot 1) = 5$: three shorter words
($\varepsilon, a, b$) plus two smaller words of length two
($aa$, $ab$), in agreement with the listing. Reversing the declared
ranking ($\rho(b) = 0$, $\rho(a) = 1$) reverses every lexicographic
tiebreak: the enumeration becomes $\varepsilon, b, a, bb, ba, ab, aa,
\dots$, and a denotation reachable by both $ab$ and $ba$ (should the
generators commute on their carrier) changes its canonical
representative from $ab$ to $ba$. Nothing mathematical broke ---
both enumerations are lawful well-orders --- but the \emph{selection}
changed, which is why the kernel treats the catalog order as declared,
versioned data (Remark~\ref{ch02:rem:slex-typed} below).
\end{example}

\begin{remark}[Inheritance by typed words; dependence on the ranking]\label{ch02:rem:slex-typed}
Both Proposition~\ref{ch02:prop:slex-total} and
Theorem~\ref{ch02:thm:slex-wellorder} are statements about arbitrary
subsets of $G^{*}$, so they hold verbatim for the kernel's typed words,
which form such a subset once start objects are fixed (formally, one takes
a disjoint union over start objects, ordering within each; Chapter~8 makes
this precise). Note also what the results do \emph{not} give: canonicity
is relative to the ranking $\rho$. Two catalogs with the same generators
in different declared orders induce different shortlex orders and may
select different representatives for the same denotation. This is a
feature the kernel documents rather than a defect it hides ---
representatives are canonical \emph{relative to a declared, versioned
catalog}, and comparing normal forms across catalogs without a catalog
identity is meaningless. The determinism discipline of
Chapter~\ref{ch03:chap} exists precisely so that $\rho$ is always explicit
data, never an accident of container iteration.
\end{remark}

\section{Summary of the toolkit}\label{ch02:sec:summary}

For ease of later reference, the primer's yield is:

\begin{itemize}
\item \emph{Extensionality and the tuple representation}
(Section~\ref{ch02:sec:finsets}): finite total functions are mapping
tuples; equality is componentwise; composition and identity satisfy the
category laws pointwise (Proposition~\ref{ch02:prop:assoc}).
\item \emph{Fibers} (Section~\ref{ch02:sec:fibers}): nonempty fibers
partition the domain (Lemma~\ref{ch02:lem:fiber-partition}); injectivity,
surjectivity, and bijectivity are fiber-size conditions; on equal finite
cardinalities the three notions coincide
(Proposition~\ref{ch02:prop:endo-trichotomy}).
\item \emph{Counting} (Section~\ref{ch02:sec:counting}): there are
$|Y|^{|X|}$ total maps (Theorem~\ref{ch02:thm:count-maps}); the test
families' cardinalities $19{,}683$, $289$, and the falling factorials are
exact (Corollary~\ref{ch02:cor:test-counts}).
\item \emph{The ladder} (Section~\ref{ch02:sec:ladder}): monoid = one-object
category; semigroup = one-object semigroupoid; powers live in
$\mathrm{End}(X)$ and nowhere else
(Proposition~\ref{ch02:prop:ladder}, Remark~\ref{ch02:rem:where-powers-live});
all parenthesizations of a composable sequence agree
(Theorem~\ref{ch02:thm:gen-assoc}); functors preserve identities,
composites, and powers (Definition~\ref{ch02:def:functor}).
\item \emph{Functional graphs} (Section~\ref{ch02:sec:funcgraphs}): orbits
are rho-shaped with tail $\leq n - 1$ and cycle $\leq n$
(Lemma~\ref{ch02:lem:rho}); powers of $f$ have an index $\leq n - 1$ and a
period equal to the lcm of cycle lengths, dividing $n!$
(Theorem~\ref{ch02:thm:index-period}); $f^{n!}$ is idempotent
(Theorem~\ref{ch02:thm:factorial-idempotent}).
\item \emph{Hilbert spaces} (Section~\ref{ch02:sec:hilbert}): adjoints,
operator norm, PSD Gram maps; $\lVert A \rVert^{2} =
\lambda_{\max}(A^{*}A)$ (Proposition~\ref{ch02:prop:norm-eig}); the
all-ones matrix $J_m$ has spectrum $\{ m, 0^{(m-1)} \}$
(Proposition~\ref{ch02:prop:allones}); permutation similarity and direct
sums assemble spectra
(Lemmas~\ref{ch02:lem:perm-similarity}--\ref{ch02:lem:direct-sum});
projectors have norm at most one, bare idempotents do not
(Proposition~\ref{ch02:prop:projector-facts}).
\item \emph{Shortlex} (Section~\ref{ch02:sec:shortlex}): a strict total
well-order on words over a ranked alphabet
(Proposition~\ref{ch02:prop:slex-total},
Theorem~\ref{ch02:thm:slex-wellorder}); least representatives exist and
are unique per class (Corollary~\ref{ch02:cor:min-rep}), relative to the
declared ranking; the rank formula makes the enumeration computable and
auditable (Proposition~\ref{ch02:prop:slex-rank}).
\end{itemize}

Every later chapter draws on this list and on nothing outside it. The
mathematics being standard is the point: the kernel's claims are meant to
be auditable by any reader with this chapter in hand, and the program's
open formalization obligations (Chapter~11, Appendix~B) are obligations to
mechanize precisely these statements, no more and no less.
